\section{The SDF model and its discretisation}
\label{sec:sdf}

\subsection{Primitives and transformations}

A scene is a list of operations applied to an initially empty field.  Each
operation carries a distance function, a transformation, a brush, and an
optional painter.  The distance functions are evaluated in the primitive's
local space; the transformation maps the query point into that space and
the result is scaled back when the field is not unit-speed.  The engine's
primitive catalogue includes spheres, boxes, capsules, cylinders, cones,
tori, tapered primitives, height maps, road splines, swept polygons, and
extruded SVG polygons.  The exact primitive is irrelevant to the mesher:
every primitive reduces to a function
\begin{equation}
  d = f_{\text{primitive}}(\mathbf{x}_{\text{local}}),
\end{equation}
and every CSG operation reduces to a scalar combination of two such
values.

Brushes carry the material identity that eventually becomes the vertex's
\code{brushIndex}; painters write the per-node HSV tint and the mapping
mode/UV scale.  The mesher never interprets either: it copies the node's
brush and mapping state into the vertex payload in stage 1.

\subsection{Corner sampling}

The octree stores a field value only at the eight corners of each cube.
For a cube with minimum $\mathbf{c}$ and edge $L$, corner $i$ is
\begin{equation}
  C_i = \mathbf{c} + L\,\mathbf{s}_i,
  \qquad
  \mathbf{s}_i = \bigl((i\gg2)\&1,\;(i\gg1)\&1,\;i\&1\bigr),
  \label{eq:corner}
\end{equation}
with bit 2 selecting $x$, bit 1 selecting $y$, and bit 0 selecting $z$.
The stored array is
\begin{equation}
  \mathbf{d} = (d_0,\ldots,d_7), \qquad d_i = f(C_i).
\end{equation}
A corner with no data is stored as $+\infty$.  This sentinel is important:
a value of $+\infty$ means ``unknown'', not ``very far outside''.  The
classifier skips it, and the QEF skips any edge with a non-finite
endpoint.  Treating it as a finite positive value would place vertices at
the corner with a wild normal, creating the familiar ``spike through the
terrain'' artefact.

\subsection{Classification}

The node type is computed from the signs of the finite corners:
\begin{equation}
  \operatorname{type}(\mathbf{d}) =
  \begin{cases}
    \text{Empty}, & \text{no finite negative corner},\\
    \text{Solid}, & \text{finite corners are all negative},\\
    \text{Surface}, & \text{both signs are present}.
  \end{cases}
  \label{eq:type}
\end{equation}
The implementation is careful to skip $+\infty$ and to treat $0$ as
non-negative.  A cell with a negative corner and an $+\infty$ corner is
classified \emph{Solid} by the stored type yet can still emit a crossing
edge; therefore the GPU scan list is filtered by an explicit
``straddles the surface'' test, not by the stored type alone.  This
distinction is one of the subtle correctness points of the flattening.

\subsection{Trilinear reconstruction and its gradient}

Inside a cube, write the local coordinate
\begin{equation}
  \mathbf{u} = \frac{\mathbf{x}-\mathbf{c}}{L}\in[0,1]^3.
\end{equation}
The trilinear interpolant is
\begin{align}
  f(\mathbf{u}) &=
    (1-u_x)(1-u_y)(1-u_z)d_0
  + (1-u_x)(1-u_y)u_z d_1 \nonumber\\
  &\quad + (1-u_x)u_y(1-u_z)d_2
  + (1-u_x)u_y u_z d_3 \nonumber\\
  &\quad + u_x(1-u_y)(1-u_z)d_4
  + u_x(1-u_y)u_z d_5 \nonumber\\
  &\quad + u_x u_y(1-u_z)d_6
  + u_x u_y u_z d_7 .
\end{align}
The analytic gradient is used for normals:
\begin{align}
  \frac{\partial f}{\partial x} &=
    (1-u_y)(1-u_z)(d_4-d_0)
  + (1-u_y)u_z(d_5-d_1) \nonumber\\
  &\quad + u_y(1-u_z)(d_6-d_2)
  + u_y u_z(d_7-d_3),\\
  \frac{\partial f}{\partial y} &=
    (1-u_x)(1-u_z)(d_2-d_0)
  + u_x(1-u_z)(d_3-d_1) \nonumber\\
  &\quad + (1-u_x)u_z(d_6-d_4)
  + u_x u_z(d_7-d_5),\\
  \frac{\partial f}{\partial z} &=
    (1-u_x)(1-u_y)(d_1-d_0)
  + (1-u_x)u_y(d_3-d_2) \nonumber\\
  &\quad + u_x(1-u_y)(d_5-d_4)
  + u_x u_y(d_7-d_6).
\end{align}
The world-space gradient is $\nabla_{\mathbf{x}} f =
(\partial_x f,\partial_y f,\partial_z f)/L$.  The normal is the normalized
gradient.  If the gradient is degenerate, the shader falls back to
$(0,1,0)$; this is preferable to a NaN vertex.

\subsection{The trilinear lerp with sentinels}

A naive \code{mix(a,b,t)} evaluates $a(1-t)+bt$.  If one endpoint is
$+\infty$ and the corresponding weight is exactly zero, the product is
$0\cdot\infty=\mathrm{NaN}$.  The engine uses a guarded lerp:
\begin{equation}
  \operatorname{lerp}(a,b,t) =
  \begin{cases}
    a, & t\le 0,\\
    b, & t\ge 1,\\
    +\infty, & a=+\infty \text{ or } b=+\infty,\\
    a+(b-a)t, & \text{otherwise}.
  \end{cases}
  \label{eq:lerp}
\end{equation}
This keeps the sentinel semantics through interpolation and is used by the
CPU reference, the GPU shader, and the simplifier.  It is also why the
shader marks the sensitive expressions \code{precise}: reassociation can
change the sign of an expression that is exactly zero at a corner.

\subsection{The QEF vertex}

For each finite sign-changing edge $e=(i,j)$, let
\begin{equation}
  \mathbf{p}_e = (1-t_e)\mathbf{p}_i + t_e\mathbf{p}_j,
  \qquad
  t_e = \clip\!\left(\frac{d_i}{d_i-d_j},0,1\right).
\end{equation}
The edge normal is the interpolation of the two corner normals,
normalized.  The QEF is
\begin{equation}
  E(\mathbf{x}) = \sum_e \left(\mathbf{n}_e\cdot
  (\mathbf{x}-\mathbf{p}_e)\right)^2 .
\end{equation}
Its normal equations are
\begin{equation}
  \Bigl(\sum_e \mathbf{n}_e\mathbf{n}_e^{T}\Bigr)\mathbf{x}
  = \sum_e \mathbf{n}_e\,(\mathbf{n}_e\cdot\mathbf{p}_e),
  \label{eq:qef}
\end{equation}
solved by Cramer's rule.  Two regularizers are essential:
\begin{enumerate}
\item the matrix is shifted by $\epsilon I$ with $\epsilon=10^{-2}$;
\item the right-hand side is shifted by $\epsilon\,\bar{\mathbf{p}}$
  where $\bar{\mathbf{p}}$ is the average edge crossing.
\end{enumerate}
Without them, a cell with one crossing edge has a rank-one matrix and the
solution can shoot arbitrarily far away.  The result is clamped to the
cell cube, so a dual vertex can never leave its cell.

\begin{figure}[htbp]
\centering
\begin{tikzpicture}[scale=1.25]
  \draw[grid] (0,0) rectangle (3,3);
  \foreach \x in {0,1,2,3} \foreach \y in {0,1,2,3}
    \fill[ink!60] (\x,\y) circle (0.6pt);
  % a slice
  \draw[accent,thick] (0.2,2.35) .. controls (1.0,1.9) and (2.1,1.7) .. (2.8,1.2);
  \fill[accent2!12] (0.2,2.35) -- (2.8,1.2) -- (2.8,3) -- (0.2,3) -- cycle;
  % crossings and normals
  \foreach \x/\y/\nx/\ny in {0.55/2.25/-0.15/0.99, 1.45/1.88/-0.15/0.99, 2.35/1.43/-0.15/0.99} {
    \fill[accent2] (\x,\y) circle (2pt);
    \draw[accent2,-{Latex[length=1.8mm]}] (\x,\y) -- ++(\nx*0.55,\ny*0.55);
  }
  \node[accent2,font=\small] at (2.55,2.55) {$d>0$};
  \node[accent,font=\small] at (0.55,1.25) {$d<0$};
  \draw[decorate,decoration={brace,amplitude=4pt},ink!70]
    (3.12,1.43) -- (3.12,2.25) node[midway,right=5pt,font=\scriptsize]{edge crossings};
\end{tikzpicture}
\caption{A surface cell in two dimensions.  The shaded region is the
inside ($d<0$).  Sign changes on cube edges give crossing points and
normals; the QEF fits a dual vertex to those planes.  In three dimensions
there are twelve cube edges, and the vertex is clamped inside the cube.}
\label{fig:qef2d}
\end{figure}

\subsection{Brushes, painters, and mapping state}

The node stores the brush index, an HSV tint, the texture-mapping mode, and
two generic data values.  The brush is painted at the shape leaf; internal
nodes inherit the dominant painted child.  The GPU vertex payload packs:
\begin{center}
\begin{tabular}{@{}lll@{}}
\toprule
Word & Components & Source\\
\midrule
0 & position.xyz, color.x & QEF position; color is 1\\
1 & color.yz, texCoord.xy & texCoord = node data (UV scale)\\
2 & normal.xyz, brushIndex bits & normalized SDF gradient; brush\\
3 & textureType bits, hsv & mapping mode; HSV tint\\
\bottomrule
\end{tabular}
\end{center}
The renderer treats the 64-byte record as the canonical \code{Vertex}; the
same record is used by the raster path, the ray-tracing BLAS, and the
offline exporter.  This is why the finalize stage does not need a
conversion pass.
